\subsection{Lavjerrësi jolinear dhe perioda e varur nga amplituda}
Ekuacioni pa fërkim është
\begin{equation}
 \ddot\theta+\frac{g}{L}\sin\theta=0.
\end{equation}
Linearizimi $\sin\theta\simeq\theta$ jep periodën $T_0=2\pi\sqrt{L/g}$, por modeli linear humbet varësinë nga amplituda. Nga ruajtja e energjisë,
\begin{equation}
 \frac12L^2\dot\theta^2+gL(1-\cos\theta)
 =gL(1-\cos\theta_0),
\end{equation}
merret perioda
\begin{equation}
 T=4\sqrt{\frac{L}{g}}
 \int_0^{\pi/2}\frac{\dd\phi}
 {\sqrt{1-k^2\sin^2\phi}},
 \qquad k=\sin\frac{\theta_0}{2}.
\end{equation}
Ky integral eliptik mund të llogaritet me kuadraturë dhe të krahasohet me periodën e nxjerrë nga integrimi kohor. Për amplituda të vogla,
\begin{equation}
 \frac{T}{T_0}=1+\frac{\theta_0^2}{16}
 +\frac{11\theta_0^4}{3072}+\cdots.
\end{equation}
Krahasimi i tri rrugëve --- seria, kuadratura dhe ODE-ja --- është verifikim i fuqishëm ndërmetodik.

Perioda numerike nga trajektorja matet me kalime në zero ose me kulme. Interpolimi ndërmjet dy hapave zvogëlon gabimin e kuantizimit kohor. Për integrim afatgjatë, monitorohet energjia. Një metodë jo simplektike mund të japë periodë të mirë për disa cikle dhe pastaj të krijojë drift që ndryshon artificialisht amplitudën.

\subsubsection{Problemi i dy trupave dhe reduktimi në një trup}
Për masa $m_1,m_2$ me pozicione $\vect r_1,\vect r_2$, forcat janë të barabarta e të kundërta. Qendra e masës
\begin{equation}
 \vect R=\frac{m_1\vect r_1+m_2\vect r_2}{m_1+m_2}
\end{equation}
lëviz në mënyrë uniforme, ndërsa koordinata relative $\vect r=\vect r_1-\vect r_2$ plotëson
\begin{equation}
 \mu\ddot{\vect r}=-\frac{Gm_1m_2}{r^3}\vect r,
 \qquad \mu=\frac{m_1m_2}{m_1+m_2}.
\end{equation}
Pas pjesëtimit me $\mu$ merret $\ddot{\vect r}=-G(m_1+m_2)\vect r/r^3$. Invariantet janë energjia
\begin{equation}
 E=\frac12\mu|\dot{\vect r}|^2-\frac{Gm_1m_2}{r}
\end{equation}
dhe momenti këndor $\vect L=\mu\vect r\times\dot{\vect r}$. Për orbitë të lidhur me gjysmëbosht të madh $a$,
\begin{equation}
 T^2=\frac{4\pi^2}{G(m_1+m_2)}a^3.
\end{equation}
Një laborator i plotë nxjerr periodën nga simulimi për disa $a$, përshtat pjerrësinë në shkallë logaritmike dhe kontrollon konstanten fizike, jo vetëm eksponentin $3/2$.

\subsubsection{Gabimi në ngjarje dhe hapa adaptivë}
Shumë probleme kërkojnë kohën e një ngjarjeje: goditja e tokës, kalimi në periapsis ose arritja e një pragu. Nëse ngjarja ndodh ndërmjet $t_n$ dhe $t_{n+1}$, kthimi i thjeshtë i $t_{n+1}$ ka gabim të rendit $h$. Duhet interpoluar funksioni i ngjarjes $g(t,y)$ dhe zgjidhur $g=0$ brenda hapit. Bibliotekat ODE e bëjnë këtë me interpolues të dendur, por përdoruesi duhet të përcaktojë drejtimin e kalimit dhe nëse integrimi ndalet.

Kontrolli adaptiv përdor dy vlerësime me rende të ndryshme ose step-doubling. Nëse gabimi lokal i vlerësuar është $e$, hapi i ri zgjidhet afërsisht
\begin{equation}
 h_{\mathrm{ri}}=s h\left(\frac{\mathrm{tol}}{e}\right)^{1/(p+1)},
\end{equation}
me faktor sigurie $s<1$. Toleranca absolute mbron komponentët pranë zeros; toleranca relative shkallëzon komponentët e mëdhenj. Në sisteme me njësi të ndryshme nevojiten toleranca komponentore ose nondimensionalizim.
